% LSE BSc Dissertation in Mathematics
\documentclass[a4paper,oneside,11pt]{report}

\input{template-preamble}

\usepackage{newpxtext,newpxmath}
\usepackage{inconsolata}
\usepackage{roboto}
\usepackage{microtype}

% Additional packages for this dissertation
\usepackage{algorithm2e}
\usepackage{mathtools}
\usepackage{dsfont} % for indicator function \mathds{1}
\usepackage{booktabs} % for \toprule, \midrule, \bottomrule

% Custom commands
\newcommand{\ind}{\mathds{1}}
\newcommand{\Prob}{\mathbb{P}}
\newcommand{\R}{\mathbb{R}}
\newcommand{\N}{\mathbb{N}}
\newcommand{\calS}{\mathcal{S}}
\newcommand{\calN}{\mathcal{N}}
\newcommand{\calA}{\mathcal{A}}
\newcommand{\calF}{\mathcal{F}}
\newcommand{\calB}{\mathcal{B}}
\newcommand{\calE}{\mathcal{E}}
\newcommand{\calG}{\mathcal{G}}
\newcommand{\calU}{\mathcal{U}}
\newcommand{\calQ}{\mathcal{Q}}
\newcommand{\ip}[2]{\langle #1,\, #2 \rangle}
\newcommand{\norm}[1]{\| #1 \|}
\newcommand{\projPi}{\operatorname{proj}_\Pi}
\newcommand{\diam}{\operatorname{diam}}
\DeclareMathOperator*{\argmin}{arg\,min}
\DeclareMathOperator*{\argmax}{arg\,max}

% Theorem numbering within chapters
\numberwithin{theorem}{chapter}
\numberwithin{equation}{chapter}

\dissertationtitle{Convergence of Policy Gradient Methods\\
  to Nash Equilibria in Stochastic Games}

\candnumber{60103}

\parskip1ex

\begin{document}

\dissertationtitlepage

\begin{dissertationsummary}
We give a self-contained, measure-theoretic treatment of the convergence
of policy gradient methods to Nash equilibria in general stochastic games,
following the recent work of Giannou, Lotidis, Mertikopoulos, and
Vlatakis-Gkaragkounis~\cite{giannou2022}. Their main result establishes
that, for any stable Nash policy in an episodic stochastic game, the
sequence of policies generated by the policy gradient algorithm converges
to the equilibrium with probability at least $1-\delta$, provided the
initial policy lies in a suitable neighbourhood and the step-size is
sufficiently small. The proof rests on a stochastic Lyapunov argument
that combines an energy inequality, martingale stability estimates, and
the Robbins--Siegmund almost supermartingale convergence lemma.

Our first contribution is expository: we develop the full proof of this
convergence theorem with all measure-theoretic foundations made explicit.
We carefully construct the underlying probability space and filtration,
verify the martingale and submartingale properties of the error
processes, and state precisely the classical results from martingale
theory that are invoked at each step.

Our second contribution is original: we prove that a simple random
restart strategy achieves global, almost-sure convergence to a Nash
equilibrium, removing the dependence on initialisation. The expected
number of restarts is bounded by the inverse of the basin probability,
and the framework naturally enables Pareto-optimal equilibrium
selection. We validate both results empirically on normal-form and
stochastic games, demonstrating that state-contingent policies
strictly outperform state-blind reductions when $|\calS| > 1$.
\end{dissertationsummary}

\dissertationtableofcontents

%% ================================================================
%%  CHAPTER 1: INTRODUCTION
%% ================================================================
\chapter{Introduction}
\label{ch:intro}

\section{Context and motivation}

A stochastic game is a multi-agent extension of the Markov decision
process (MDP), introduced by Shapley~\cite{shapley1953}. At each stage
of the game, the players find themselves in a common state, choose
actions simultaneously, receive individual rewards, and transition to a
new state according to a probability kernel that depends on the current
state and the joint action profile. The players' objectives are coupled:
each player's reward depends not only on her own action but on the
actions of all other players.

The central solution concept in this setting is the \emph{Nash
  equilibrium}~\cite{nash1951}: a joint policy from which no player can
profitably deviate. The existence of Nash equilibria in finite stochastic
games was established by Fink~\cite{fink1964}, but the question of
whether boundedly rational agents can \emph{learn} to play a Nash
equilibrium through repeated interaction remains largely open.

A natural approach is \emph{policy gradient} (PG) methods: each player
maintains a parameterised policy, plays the game for one episode,
observes a gradient signal, and updates her policy by a gradient
step~\cite{sutton2018,agarwal2021}. Despite the success of policy
gradient methods in single-agent reinforcement learning, their
convergence properties in multi-agent settings are poorly understood
outside of special classes of games such as potential games or two-player
zero-sum games~\cite{leonardos2022,zhang2022,daskalakis2022}.

\section{The main result}

The paper of Giannou, Lotidis, Mertikopoulos, and
Vlatakis-Gkaragkounis~\cite{giannou2022} addresses this gap. Their
Theorem~1 establishes that, in a \emph{general} finite stochastic game,
the sequence of policies generated by the policy gradient algorithm
converges to any \emph{stable} Nash policy with probability at least
$1-\delta$, for any $\delta > 0$, provided the initial policy lies in a
neighbourhood of the equilibrium and the step-size is chosen
appropriately. This is a local convergence guarantee: it holds
conditionally on initialisation near the equilibrium, but it applies to
general stochastic games without structural assumptions such as
potentiality or zero-sum structure.

The result is significant for several reasons. It concerns the
\emph{actual} sequence of play, not a time-averaged or best-iterate
variant. It accommodates stochastic gradient estimates with unbounded
variance, as arise from the REINFORCE estimator~\cite{williams1992}. And
it provides the first trajectory convergence guarantee for policy
gradient methods in this level of generality.

\section{Contribution and outline}

This dissertation presents a self-contained, measure-theoretic
exposition of the proof of Theorem~1. The original proof
in~\cite{giannou2022} is spread across several appendices and relies on
a number of probabilistic arguments whose measure-theoretic foundations
are left to the reader. Our aim is to make every step fully rigorous and
explicit.

The exposition is organised as follows. Chapter~\ref{ch:prelim}
introduces the model: stochastic games, policies, value functions,
solution concepts, and the policy gradient algorithm. We also set up the
probability-theoretic framework---the underlying probability space,
filtration, and the key results from martingale theory that will be
needed. Chapter~\ref{ch:proof} contains the proof of the main theorem,
developed through a sequence of intermediate results: an energy
inequality, error aggregation via martingales and submartingales,
stochastic stability, extraction of a convergent subsequence, and
convergence of the energy function. Chapter~\ref{ch:global} addresses
the natural follow-up: can we remove the dependence on initialisation?
We prove that a restart strategy converges to a Nash equilibrium almost
surely and bound the expected number of restarts.
Chapter~\ref{ch:conclusion} offers concluding remarks and directions
for future research.

\emph{A note on attribution.} The mathematical content of Chapters~\ref{ch:prelim}
and~\ref{ch:proof} follows~\cite{giannou2022} closely.
Chapter~\ref{ch:global} is our original contribution: it uses
Theorem~\ref{thm:main} as a black box to derive a global convergence
guarantee via a restart argument. The measure-theoretic treatment
throughout is ours.


%% ================================================================
%%  CHAPTER 2: PRELIMINARIES
%% ================================================================
\chapter{Stochastic Games and Policy Gradient Methods}
\label{ch:prelim}

The material in this chapter follows the setup
of~\cite{giannou2022}, to which we refer the reader for further
details and references. Throughout, we use UK English conventions.

\section{The stochastic game}
\label{sec:game}

We consider $N$-player stochastic games in the tabular, random-stopping
framework. Formally, a stochastic game is specified by a tuple
\begin{equation}
  \label{eq:game-tuple}
  \calG = \bigl(\calS,\; \calN,\; \{\calA_i, R_i\}_{i \in \calN},\;
  P,\; \zeta,\; \rho\bigr),
\end{equation}
whose components are defined as follows.
\begin{itemize}
\item $\calN = \{1, 2, \ldots, N\}$ is a finite set of \emph{agents}.
\item $\calS = \{1, 2, \ldots, S\}$ is a finite set of \emph{states}.
\item For each agent $i \in \calN$, the set
  $\calA_i = \{1, \ldots, A_i\}$ is a finite set of \emph{actions}. We
  write $\calA = \prod_{i \in \calN} \calA_i$ for the joint action
  space, and $\calA_{-i} = \prod_{j \neq i} \calA_j$ for the action
  space of all agents other than~$i$.
\item For each $i \in \calN$, the map $R_i \colon \calS \times \calA \to [-1,1]$
  is the \emph{reward function} of agent~$i$.
\item $P(s' \mid s, \alpha)$ denotes the probability of transitioning
  from state $s$ to state $s'$ when the joint action profile is
  $\alpha \in \calA$.
\item The game terminates at each state--action pair $(s,\alpha)$ with
  probability
  $\zeta_{s,\alpha} = 1 - \sum_{s' \in \calS} P(s' \mid s, \alpha) > 0$.
  We write $\zeta \coloneqq \min_{s,\alpha} \zeta_{s,\alpha}$ for the
  minimum termination probability.
\item $\rho \in \Delta(\calS)$ is the distribution of the initial state.
\end{itemize}

\paragraph{Episodic setting.}
Play proceeds in \emph{episodes}. In each episode, the initial state
$s_0$ is drawn from $\rho$. At each time step $t \ge 0$, all agents
observe the common state $s_t$, select actions $\alpha_t \in \calA$, and
receive rewards $\{R_i(s_t, \alpha_t)\}_{i \in \calN}$. The game then
terminates with probability $\zeta_{s_t, \alpha_t}$, or transitions to
$s_{t+1}$ drawn from $P(\cdot \mid s_t, \alpha_t)$. We denote the
resulting trajectory by
$\tau = (s_t, \alpha_t, r_t)_{t \le T(\tau)}$, where $T(\tau)$ is the
(random) termination time.

\section{Policies and value functions}
\label{sec:policies}

We consider \emph{stationary Markovian} policies: maps from states to
distributions over actions that do not depend on time or history. A
\emph{policy} for agent $i$ is a map
$\pi_i \colon \calS \to \Delta(\calA_i)$, which we view as an element
of the set $\Pi_i \coloneqq \Delta(\calA_i)^S$. A \emph{policy profile}
is a tuple $\pi = (\pi_i)_{i \in \calN} \in \Pi \coloneqq \prod_{i \in \calN} \Pi_i$.
We write $\lambda_\Pi$ for the Lebesgue measure on $\Pi$ (the product
of the Lebesgue measures on each simplex $\Pi_i$).

The \emph{value function} of agent $i$ under policy profile $\pi$,
starting from state $s$, is
\begin{equation}
  \label{eq:value-function}
  V_{i,s}(\pi) \coloneqq \E_{\tau \sim \mathrm{MDP}}\biggl[
    \sum_{t=0}^{T(\tau)} R_i(s_t, \alpha_t) \biggm| s_0 = s
  \biggr],
\end{equation}
where the expectation is over the randomness induced by $\pi$ and the
transition kernel~$P$. Setting
$V_{i,\rho}(\pi) \coloneqq \E_{s \sim \rho}[V_{i,s}(\pi)]$, we obtain
the \emph{expected value} under initial distribution $\rho$.

\paragraph{Discounted state visitation distribution.}
For a policy profile $\pi \in \Pi$ and initial distribution $\rho$,
define the \emph{unnormalised discounted state visitation measure}
\begin{equation}
  \label{eq:visitation}
  \bar{d}^\pi_\rho(s) = \E_{\tau \sim \mathrm{MDP}}\biggl[
    \sum_{t=0}^{T(\tau)} \ind\{s_t = s\} \biggm| s_0 \sim \rho
  \biggr],
\end{equation}
and the normalised version
$d^\pi_\rho(s) = \bar{d}^\pi_\rho(s) / Z^\pi_\rho$, where
$Z^\pi_\rho = \sum_{s \in \calS} \bar{d}^\pi_\rho(s)$. Since the
termination probability satisfies $\zeta > 0$, we have
$\bar{d}^\pi_\rho(s) \le 1/\zeta < \infty$ for all $s$, so these
quantities are well-defined.

\paragraph{The mismatch coefficient.}
The quantity
\begin{equation}
  \label{eq:mismatch}
  C_\calG \coloneqq \max_{\pi, \pi' \in \Pi}
  \Bigl\{ \bigl\| \bar{d}^\pi_\rho / \bar{d}^{\pi'}_\rho
  \bigr\|_\infty \Bigr\}
\end{equation}
measures the worst-case ratio of visitation distributions across policy
profiles. Following~\cite{agarwal2021, giannou2022}, we assume
$C_\calG < \infty$, which is equivalent to
$d^\pi_\rho(s) > 0$ for all $\pi$ and $s$.

\section{Nash policies and their refinements}
\label{sec:nash}

\begin{definition}[Nash policy]
  \label{def:nash}
  A policy $\pi^* = (\pi^*_i)_{i \in \calN} \in \Pi$ is a \emph{Nash
    policy} for the initial distribution $\rho$ if, for every player
  $i \in \calN$ and every unilateral deviation
  $\pi_i \in \Delta(\calA_i)^S$,
  \begin{equation}
    \label{eq:NE}
    V_{i,\rho}(\pi^*_i;\, \pi^*_{-i}) \ge
    V_{i,\rho}(\pi_i;\, \pi^*_{-i}).
    \tag{NE}
  \end{equation}
\end{definition}

We write $v_i(\pi) = \nabla_i V_{i,\rho}(\pi)$ for the gradient of
player~$i$'s value function with respect to $\pi_i$ (keeping
$\pi_{-i}$ fixed), and $v(\pi) = (v_i(\pi))_{i \in \calN}$ for the
concatenated gradient vector.

A key structural property of stochastic games, established
in~\cite{giannou2022, daskalakis2022, zhang2022}, is the following.

\begin{lemma}[Gradient dominance property;
  {\cite[Lemma~2]{giannou2022}}]
  \label{lem:GDP}
  For any policy profile $\pi = (\pi_i)_{i \in \calN} \in \Pi$ and any
  unilateral deviation $\pi'_i \in \Pi_i$ of player~$i$,
  \begin{equation}
    \label{eq:GDP}
    V_{i,\rho}(\pi'_i;\, \pi_{-i}) - V_{i,\rho}(\pi_i;\, \pi_{-i})
    \le C_\calG \max_{\bar\pi_i \in \Pi_i}
    \ip{\nabla_i V_{i,\rho}(\pi)}{\bar\pi_i - \pi_i}.
    \tag{GDP}
  \end{equation}
\end{lemma}

An immediate consequence is that first-order stationary points of~$V$
are Nash~\cite[Lemma~3]{giannou2022}.

\begin{lemma}[First-order stationarity implies Nash]
  \label{lem:FOS-Nash}
  A policy $\pi^* \in \Pi$ is Nash if and only if it satisfies the
  first-order stationary condition
  \begin{equation}
    \label{eq:FOS}
    \ip{v(\pi^*)}{\pi - \pi^*} \le 0
    \quad \text{for all } \pi \in \Pi.
    \tag{FOS}
  \end{equation}
\end{lemma}

We now introduce refinements of Nash policies that play a central role
in the convergence analysis.

\begin{definition}[Refinements of Nash policies;
  {\cite[Definition~2]{giannou2022}}]
  \label{def:refinements}
  Let $\pi^* \in \Pi$ be a Nash policy.
  \begin{itemize}
  \item $\pi^*$ is \emph{stable} if
    $\ip{v(\pi)}{\pi - \pi^*} < 0$ for all
    $\pi \ne \pi^*$ sufficiently close to $\pi^*$.
  \item $\pi^*$ is \emph{second-order stationary} (SOS) if
    \begin{equation}
      \label{eq:SOS}
      (\pi - \pi^*)^\top \operatorname{Jac}_v(\pi^*)(\pi - \pi^*) < 0
      \quad \text{for all } \pi \in \Pi \setminus \{\pi^*\},
      \tag{SOS}
    \end{equation}
    where
    $\operatorname{Jac}_v(\pi^*) = (\nabla_j v_i(\pi^*))_{i,j \in
      \calN}$ is the Jacobian of~$v$ at $\pi^*$.
  \item $\pi^*$ is \emph{deterministic} if $\pi^*_i(s)$ places all
    mass on a single action for every $i$ and $s$.
  \item $\pi^*$ is \emph{strict} if it is deterministic and
    \eqref{eq:FOS} holds with strict inequality whenever
    $\pi \ne \pi^*$.
  \end{itemize}
\end{definition}

One has the implications
strict/deterministic $\Longrightarrow$ SOS $\Longrightarrow$ stable
$\Longrightarrow$ FOS $=$ Nash; see~\cite[Proposition~1]{giannou2022}.
The following quantitative versions of the stability condition will be
used in the proof.

\begin{proposition}[Drift inequalities;
  {\cite[Proposition~1]{giannou2022}}]
  \label{prop:drift}
  Let $\pi^* \in \Pi$ be a Nash policy.
  \begin{enumerate}[label=\textup{(\alph*)}]
  \item If $\pi^*$ is second-order stationary, there exists $\mu > 0$
    such that
    \begin{equation}
      \label{eq:drift-SOS}
      \ip{v(\pi)}{\pi - \pi^*} \le -\mu\, \norm{\pi - \pi^*}^2
      \quad \text{for all } \pi \text{ sufficiently close to } \pi^*.
    \end{equation}
  \item If $\pi^*$ is strict, there exists $\mu > 0$ such that
    \begin{equation}
      \label{eq:drift-strict}
      \ip{v(\pi)}{\pi - \pi^*} \le -\mu\, \norm{\pi - \pi^*}
      \quad \text{for all } \pi \text{ sufficiently close to } \pi^*.
    \end{equation}
  \end{enumerate}
\end{proposition}

\section{The policy gradient algorithm}
\label{sec:PG}

We now describe the inter-episode learning framework. At the beginning
of each episode $n = 1, 2, \ldots$, each agent $i$ commits to a policy
$\pi_{i,n} \in \Pi_i$. The agents play the game under the joint policy
$\pi_n = (\pi_{i,n})_{i \in \calN}$, observe a gradient signal
$\hat v_n$, and update their policies according to the \emph{policy
  gradient} (PG) template:
\begin{equation}
  \label{eq:PG}
  \pi_{n+1} = \projPi(\pi_n + \gamma_n \hat v_n),
  \tag{PG}
\end{equation}
where $\projPi \colon \prod_i \R^{A_i \times S} \to \Pi$ denotes the
Euclidean projection onto the feasible policy space, and
$\gamma_n > 0$ is the step-size. We take the step-size to be of the
form
\begin{equation}
  \label{eq:step-size}
  \gamma_n = \frac{\gamma}{(n+m)^p},
\end{equation}
for some $\gamma > 0$, $m \ge 0$, and $p \in (1/2, 1]$.

\paragraph{Gradient signal decomposition.}
The gradient signal $\hat v_n$ is decomposed as
\begin{equation}
  \label{eq:oracle-decomp}
  \hat v_n = v(\pi_n) + U_n + b_n,
\end{equation}
where
\begin{equation}
  \label{eq:U-b-def}
  U_n \coloneqq \hat v_n - \E[\hat v_n \mid \calF_n],
  \qquad
  b_n \coloneqq \E[\hat v_n \mid \calF_n] - v(\pi_n).
\end{equation}
Here $\calF_n$ is the natural filtration defined in
Section~\ref{sec:prob-framework} below. The term $U_n$ is a zero-mean
noise (a martingale difference) and $b_n$ is a systematic bias. We
assume deterministic upper bounds:
\begin{equation}
  \label{eq:bias-var-bounds}
  \E\bigl[\norm{b_n} \bigm| \calF_n\bigr] \le B_n,
  \qquad
  \E\bigl[\norm{U_n}^2 \bigm| \calF_n\bigr] \le \sigma_n^2,
\end{equation}
where $B_n = \mathcal{O}(1/n^{\ell_b})$ and
$\sigma_n = \mathcal{O}(n^{\ell_\sigma})$ for some
$\ell_b, \ell_\sigma > 0$ depending on the information model.

\paragraph{REINFORCE estimator.}
A concrete instance of this framework arises from the REINFORCE
algorithm~\cite{williams1992}. Each agent plays the game with a
\emph{sampling policy} $\tilde \pi_{i,n}$ (possibly different from
$\pi_{i,n}$) and constructs the gradient estimate
\begin{equation}
  \label{eq:REINFORCE}
  \hat v_i = R_i(\tau) \cdot \sum_{t=0}^{T(\tau)}
  \nabla_i\bigl(\log \tilde\pi_i(\alpha_{i,t} \mid s_t)\bigr).
\end{equation}
In the $\varepsilon$-Greedy PG variant, one sets
$\tilde\pi_{i,n} = (1-\varepsilon_n)\pi_{i,n} + \varepsilon_n \,
\mathrm{Unif}_{\calA_i}$. This introduces bias of order
$B_n = \mathcal{O}(\varepsilon_n)$ and bounds the variance at
$\sigma_n^2 = \mathcal{O}(1/\varepsilon_n)$;
see~\cite[Lemma~4]{giannou2022}.

\section{Probability-theoretic framework}
\label{sec:prob-framework}

We now set up the measure-theoretic foundations that underpin the
convergence analysis. This section is not present
in~\cite{giannou2022}; it makes explicit the structure that is used
implicitly throughout their proofs.

\paragraph{The probability space.}
We work on a complete probability space $(\Omega, \calF, \Prob)$ rich
enough to support the entire learning process: the initial policy
$\pi_1$, all episode trajectories $\tau_1, \tau_2, \ldots$, and hence
all subsequent policies $\pi_2, \pi_3, \ldots$ determined
by~\eqref{eq:PG}. Formally, one may construct $(\Omega, \calF, \Prob)$
as the product of the law of~$\pi_1$ with a countable product of copies
of the trajectory space, using the Ionescu-Tulcea extension
theorem~\cite{durrett2019}; the precise construction is standard and we
omit the details.

\paragraph{Filtration.}
The natural filtration for this process is
\begin{equation}
  \label{eq:filtration}
  \calF_n \coloneqq \sigma\bigl(\pi_1,\, \hat v_1,\, \hat v_2,\,
  \ldots,\, \hat v_{n-1}\bigr)
  \quad \text{for } n \ge 1,
\end{equation}
with $\calF_1 = \sigma(\pi_1)$. This $\sigma$-algebra contains all
information available \emph{before} episode $n$ is played. In
particular:
\begin{itemize}
\item $\pi_n$ is $\calF_n$-measurable, since $\pi_n$ is a deterministic
  function of $\pi_{n-1}$ and $\hat v_{n-1}$ via~\eqref{eq:PG}.
\item $\hat v_n$ is \emph{not} $\calF_n$-measurable: it depends on the
  trajectory $\tau_n$, which is generated \emph{during} episode $n$.
\item $\hat v_n$ is $\calF_{n+1}$-measurable by definition.
\end{itemize}

Since $\pi_n$ is $\calF_n$-measurable and $v$ is a deterministic
function, $v(\pi_n)$ is $\calF_n$-measurable. The bias
$b_n = \E[\hat v_n \mid \calF_n] - v(\pi_n)$ is likewise
$\calF_n$-measurable. The noise $U_n$ and all quantities derived from
it (such as $\xi_n$, defined in Chapter~\ref{ch:proof}) are
$\calF_{n+1}$-measurable.

\paragraph{Conditional independence between episodes.}
The key structural property is that each episode trajectory $\tau_n$ is
conditionally independent of all previous trajectories given the
policy~$\pi_n$. More precisely, for any bounded measurable function
$f$,
\begin{equation}
  \label{eq:cond-ind}
  \E\bigl[f(\tau_n) \bigm| \calF_n\bigr]
  = \E\bigl[f(\tau_n) \bigm| \pi_n\bigr]
  \quad \text{a.s.}
\end{equation}
This follows from the Markov property of the episodic game: given
$\pi_n$, the distribution of $\tau_n$ is completely determined by the
MDP dynamics and does not depend on previous trajectories. It is this
property that ensures $\E[U_n \mid \calF_n] = 0$.

\paragraph{Classical results from martingale theory.}
We shall invoke the following standard results. We state them here for
easy reference; proofs can be found
in~\cite{williams1991,durrett2019,hall1980}.

\begin{theorem}[Doob's maximal inequality;
  {\cite[Theorem~4.4.2]{durrett2019}}]
  \label{thm:doob-max}
  Let $(X_n)_{n \ge 0}$ be a nonnegative submartingale. Then for any
  $\lambda > 0$,
  \begin{equation}
    \Prob\Bigl(\max_{0 \le k \le n} X_k \ge \lambda\Bigr)
    \le \frac{\E[X_n]}{\lambda}.
  \end{equation}
\end{theorem}

\begin{theorem}[Strong law for martingale differences;
  {\cite[Theorem~2.18]{hall1980}}]
  \label{thm:SLLN-mart}
  Let $(M_n)_{n \ge 1}$ be a square-integrable martingale with
  differences $d_n = M_n - M_{n-1}$. If $c_n \to \infty$ and
  \begin{equation}
    \sum_{n=1}^\infty \frac{\E[d_n^2]}{c_n^2} < \infty,
  \end{equation}
  then $M_n / c_n \to 0$ almost surely.
\end{theorem}

\begin{theorem}[Doob's submartingale convergence;
  {\cite[Theorem~4.2.11]{durrett2019}}]
  \label{thm:doob-conv}
  If $(X_n)_{n \ge 1}$ is a submartingale satisfying
  $\sup_n \E[X_n^+] < \infty$, then $X_n$ converges almost surely to a
  finite limit.
\end{theorem}

\begin{theorem}[Robbins--Siegmund lemma;
  {\cite{robbins1971}, \cite[p.~49]{polyak2020}}]
  \label{thm:robbins-siegmund}
  Let $(\calF_n)$ be a filtration and let $(X_n)$, $(a_n)$, $(b_n)$ be
  nonnegative adapted sequences satisfying
  \begin{equation}
    \E[X_{n+1} \mid \calF_n] \le X_n + a_n - b_n
    \quad \text{a.s.\ for all } n,
  \end{equation}
  with $\sum_{n=1}^\infty a_n < \infty$ a.s. Then $X_n$ converges
  almost surely to a finite limit and
  $\sum_{n=1}^\infty b_n < \infty$ a.s.
\end{theorem}

\paragraph{Notation.}
We write $\norm{\Pi} \coloneqq \max_{\pi, \pi' \in \Pi} \norm{\pi - \pi'}$
for the diameter of the policy space, and
$G \coloneqq \sup_\pi \norm{v(\pi)}$ for the gradient bound; both are
finite since $\Pi$ is compact and $v$ is
continuous~\cite[Lemmas~D.5, D.7]{giannou2022}.


%% ================================================================
%%  CHAPTER 3: THE PROOF
%% ================================================================
\chapter{Asymptotic Convergence to Stable Nash Policies}
\label{ch:proof}

This chapter contains the proof of the main convergence result. The
exposition follows Appendix~A of~\cite{giannou2022}, with all
measure-theoretic details made explicit.

\section{Statement of the main theorem}
\label{sec:thm-statement}

\begin{theorem}[Convergence to stable Nash policies;
  {\cite[Theorem~1]{giannou2022}}]
  \label{thm:main}
  Let\/ $\pi^*$ be a stable Nash policy, and let $\pi_n$ be the
  sequence of play generated by~\eqref{eq:PG} with step-size
  $\gamma_n = \gamma/(n+m)^p$, $p \in (1/2, 1]$, and policy gradient
  estimates such that $p + \ell_b > 1$ and
  $p - \ell_\sigma > 1/2$, as per~\eqref{eq:bias-var-bounds}. Then
  there exists a neighbourhood $\calU$ of\/ $\pi^*$ in $\Pi$ such that,
  for any given $\delta > 0$, we have
  \begin{equation}
    \label{eq:thm-main}
    \Prob\bigl(\pi_n \to \pi^* \mid \pi_1 \in \calU\bigr)
    \ge 1 - \delta,
  \end{equation}
  provided that $\gamma$ is small enough (or $m$ large enough) relative
  to~$\delta$.
\end{theorem}

\section{Overview of the proof strategy}
\label{sec:proof-overview}

The proof proceeds in five steps, which we summarise before carrying
them out in detail.

\begin{enumerate}
\item \textbf{Energy inequality} (Section~\ref{sec:energy-ineq}). We
  show that the squared distance
  $D(\pi) = \tfrac{1}{2}\norm{\pi - \pi^*}^2$ acts as a local Lyapunov
  function: $D(\pi_n)$ decreases on average near $\pi^*$, up to error
  terms arising from the noise and bias.

\item \textbf{Error aggregation} (Section~\ref{sec:error-agg}). We
  aggregate the error terms into a martingale $M_n$ (zero-mean errors)
  and a submartingale $S_n$ (bias and variance), and study their sum
  $W_n = M_n^2 + S_n$.

\item \textbf{Stochastic stability} (Section~\ref{sec:stability}). We
  show that, with probability at least $1-\delta$, the iterates $\pi_n$
  remain within a ball $\calB$ centred at $\pi^*$ for all $n$. The key
  tool is Doob's maximal inequality, applied to bound the probability of
  ``bad realisations'' of the noise.

\item \textbf{Convergent subsequence}
  (Section~\ref{sec:conv-subseq}). Using the strong law of large numbers
  for martingale differences and Doob's convergence theorem, we extract
  a subsequence $\pi_{n_k} \to \pi^*$ with probability one on the
  stability event.

\item \textbf{Convergence of the energy} (Section~\ref{sec:energy-conv}).
  We show that $D_n = D(\pi_n)$ converges to a finite limit using the
  Robbins--Siegmund lemma. Combined with the existence of a convergent
  subsequence, this forces $D_n \to 0$, i.e., $\pi_n \to \pi^*$.
\end{enumerate}

\section{The energy inequality}
\label{sec:energy-ineq}

Define the \emph{energy function}
\begin{equation}
  \label{eq:energy}
  D(\pi) \coloneqq \tfrac{1}{2}\norm{\pi - \pi^*}^2,
\end{equation}
and set $D_n \coloneqq D(\pi_n)$ for the energy along the iterates.

\begin{lemma}[Energy inequality; {\cite[Lemma~A.1]{giannou2022}}]
  \label{lem:energy-ineq}
  For all $n \ge 1$,
  \begin{equation}
    \label{eq:energy-ineq}
    D_{n+1} \le D_n + \gamma_n \ip{v(\pi_n)}{\pi_n - \pi^*}
    + \gamma_n \xi_n + \gamma_n \chi_n + \gamma_n^2 \psi_n^2,
  \end{equation}
  where the error terms are
  \begin{equation}
    \label{eq:error-terms}
    \xi_n = \ip{U_n}{\pi_n - \pi^*}, \qquad
    \chi_n = \norm{\Pi}\, B_n, \qquad
    \psi_n^2 = \tfrac{1}{2}\norm{\hat v_n}^2.
  \end{equation}
\end{lemma}

\begin{proof}
  By the definition of the iterates~\eqref{eq:PG} and the fact that
  $\pi^* \in \Pi$,
  \begin{align}
    D_{n+1}
    &= \tfrac{1}{2}\norm{\pi_{n+1} - \pi^*}^2
    = \tfrac{1}{2}\norm{\projPi(\pi_n + \gamma_n \hat v_n)
      - \projPi(\pi^*)}^2.
      \label{eq:D-expand-1}
  \end{align}
  Since the Euclidean projection onto a closed convex set is
  non-expansive (i.e., $\norm{\projPi(x) - \projPi(y)} \le \norm{x-y}$
  for all $x, y$), and $\projPi(\pi^*) = \pi^*$, we obtain
  \begin{equation}
    D_{n+1} \le \tfrac{1}{2}\norm{\pi_n + \gamma_n \hat v_n - \pi^*}^2.
    \label{eq:nonexpansive}
  \end{equation}
  Expanding the squared norm gives
  \begin{align}
    D_{n+1}
    &\le \tfrac{1}{2}\norm{\pi_n - \pi^*}^2
       + \gamma_n \ip{\hat v_n}{\pi_n - \pi^*}
       + \tfrac{1}{2}\gamma_n^2 \norm{\hat v_n}^2 \notag\\
    &= D_n + \gamma_n \ip{\hat v_n}{\pi_n - \pi^*}
       + \gamma_n^2 \psi_n^2.
       \label{eq:D-expanded}
  \end{align}
  Substituting the decomposition
  $\hat v_n = v(\pi_n) + U_n + b_n$ from~\eqref{eq:oracle-decomp}, the
  inner product becomes
  \begin{equation}
    \ip{\hat v_n}{\pi_n - \pi^*}
    = \ip{v(\pi_n)}{\pi_n - \pi^*}
    + \underbrace{\ip{U_n}{\pi_n - \pi^*}}_{= \xi_n}
    + \ip{b_n}{\pi_n - \pi^*}.
    \label{eq:ip-decomp}
  \end{equation}
  For the bias term, the Cauchy--Schwarz inequality yields
  \begin{equation}
    \ip{b_n}{\pi_n - \pi^*}
    \le \norm{b_n} \cdot \norm{\pi_n - \pi^*}
    \le \norm{b_n} \cdot \norm{\Pi}
    \le B_n \cdot \norm{\Pi} = \chi_n,
    \label{eq:bias-bound}
  \end{equation}
  where in the last inequality we used the conditional bound
  $\norm{b_n} \le B_n$ that holds with probability one (since
  $B_n$ is a deterministic upper bound on
  $\E[\norm{b_n} \mid \calF_n]$, and we bound $\norm{b_n}$
  deterministically by the diameter of $\Pi$; more precisely, the bound
  $\ip{b_n}{\pi_n-\pi^*} \le \norm{\Pi}\, B_n$ follows from taking the
  expectation of $\ip{b_n}{\pi_n - \pi^*}$ conditional on $\calF_n$
  and using Cauchy--Schwarz:
  $\E[\ip{b_n}{\pi_n-\pi^*} \mid \calF_n]
  \le \norm{\pi_n-\pi^*}\, \E[\norm{b_n}\mid\calF_n]
  \le \norm{\Pi}\, B_n$). Combining~\eqref{eq:D-expanded}--\eqref{eq:bias-bound}
  gives~\eqref{eq:energy-ineq}.
\end{proof}

\begin{remark}
  \label{rem:energy-roles}
  The terms in~\eqref{eq:energy-ineq} play distinct roles.  The drift
  term $\gamma_n\ip{v(\pi_n)}{\pi_n - \pi^*}$ is \emph{negative} near
  $\pi^*$ by the stability condition, driving the energy downward.  The
  noise term $\gamma_n\xi_n$ has zero conditional mean and will be
  controlled via martingale theory.  The bias term $\gamma_n\chi_n$ and
  the variance term $\gamma_n^2\psi_n^2$ are nonnegative errors that
  must be shown to be summable.
\end{remark}

\section{Error aggregation}
\label{sec:error-agg}

We now aggregate the errors appearing in~\eqref{eq:energy-ineq}.
Define, for $n \ge 1$,
\begin{equation}
  \label{eq:M-S-def}
  M_n \coloneqq \sum_{k=1}^n \gamma_k \xi_k,
  \qquad
  S_n \coloneqq \sum_{k=1}^n
  \bigl[\gamma_k \chi_k + \gamma_k^2 \psi_k^2\bigr].
\end{equation}

\begin{lemma}[Martingale and submartingale properties]
  \label{lem:mart-props}
  With respect to the filtration $(\calF_{n+1})_{n \ge 0}$:
  \begin{enumerate}[label=\textup{(\alph*)}]
  \item $(M_n)_{n \ge 1}$ is a square-integrable martingale.
  \item $(S_n)_{n \ge 1}$ is a submartingale. In fact, $S_n$ is
    nondecreasing, since each summand is nonnegative.
  \end{enumerate}
\end{lemma}

\begin{proof}
  \emph{Adaptedness.}  Since $\xi_k = \ip{U_k}{\pi_k - \pi^*}$ and
  $U_k$ is $\calF_{k+1}$-measurable (see Section~\ref{sec:prob-framework}),
  each summand $\gamma_k \xi_k$ is $\calF_{k+1}$-measurable. Thus $M_n$
  is $\calF_{n+1}$-measurable, i.e., $(M_n)$ is adapted to
  $(\calF_{n+1})_{n \ge 0}$. The same argument applies to~$S_n$.

  \emph{Integrability of~$M_n$.}  We have
  $|\xi_k| \le \norm{U_k} \cdot \norm{\Pi}$, so
  $\E[\xi_k^2] \le \norm{\Pi}^2\, \E[\norm{U_k}^2]
  \le \norm{\Pi}^2\, \sigma_k^2 < \infty$. Since
  $M_n = \sum_{k=1}^n \gamma_k \xi_k$ is a finite sum of
  square-integrable random variables, $M_n$ is square-integrable.

  \emph{Martingale property.}
  It suffices to show $\E[\gamma_n \xi_n \mid \calF_n] = 0$. Since
  $\gamma_n$ is deterministic and $\pi_n - \pi^*$ is
  $\calF_n$-measurable, we have
  \begin{equation}
    \E[\gamma_n \xi_n \mid \calF_n]
    = \gamma_n \ip{\E[U_n \mid \calF_n]}{\pi_n - \pi^*} = 0,
    \label{eq:mart-check}
  \end{equation}
  where we used $\E[U_n \mid \calF_n] = 0$ (which follows from
  $U_n = \hat v_n - \E[\hat v_n \mid \calF_n]$ by definition) and the
  fact that a $\calF_n$-measurable factor can be pulled out of a
  conditional expectation on $\calF_n$. By the tower property,
  $\E[M_n \mid \calF_n] = M_{n-1} + \E[\gamma_n\xi_n \mid \calF_n] = M_{n-1}$.

  \emph{Submartingale property of $S_n$.}  Each increment
  $\gamma_k \chi_k + \gamma_k^2 \psi_k^2 \ge 0$ (since $\chi_k, \psi_k^2$
  are nonneg), so $S_n \ge S_{n-1}$ pointwise.
\end{proof}

Following~\cite{hsieh2021,giannou2022}, we introduce the \emph{mean
  square error process}
\begin{equation}
  \label{eq:W-def}
  W_n \coloneqq M_n^2 + S_n.
\end{equation}

\begin{lemma}
  \label{lem:W-submart}
  The process $(W_n)_{n \ge 1}$ is a submartingale with respect to
  $(\calF_{n+1})$.
\end{lemma}

\begin{proof}
  Since $M_n$ is a martingale, $(M_n^2)$ is a submartingale by Jensen's
  inequality applied to the convex function $x \mapsto x^2$:
  $\E[M_n^2 \mid \calF_n] \ge (\E[M_n \mid \calF_n])^2 = M_{n-1}^2$.
  Since $S_n$ is also a submartingale, their sum $W_n = M_n^2 + S_n$ is
  a submartingale.
\end{proof}


\section{Error control and stochastic stability}
\label{sec:stability}

Let $\calB = \{\pi \in \Pi : \norm{\pi - \pi^*} \le \varrho\}$ be a
ball of radius $\varrho$ centred at $\pi^*$ in $\Pi$, chosen so that
$\ip{v(\pi)}{\pi - \pi^*} < 0$ for all
$\pi \in \calB \setminus \{\pi^*\}$ (this is possible by the stability
of $\pi^*$; without loss of generality, we take $\calB$ maximal in this
regard).

Fix a parameter $a > 0$ satisfying $2a + \sqrt{a} < \varrho$, and
define the neighbourhood
\begin{equation}
  \label{eq:U-def}
  \calU \coloneqq \{\pi \in \Pi : D(\pi) \le a\}
  = \bigl\{\pi \in \Pi : \norm{\pi - \pi^*}^2/2 \le a\bigr\}.
\end{equation}
We assume $\pi_1 \in \calU$, i.e., the process is initialised in a ball
of radius $\sqrt{2a}$ around~$\pi^*$.

Define the ``good'' events
\begin{equation}
  \label{eq:good-events}
  \calE_n \coloneqq \{\pi_k \in \calB
  \text{ for all } k = 1, \ldots, n\},
  \qquad
  H_n \coloneqq \{W_k \le a
  \text{ for all } k = 1, \ldots, n\},
\end{equation}
with the convention $\calE_0 = H_0 = \Omega$.

The following technical lemma provides the recursive estimates needed for
stability. Its proof is the most involved part of the argument.

\begin{lemma}[Error control; {\cite[Lemma~A.2]{giannou2022}}]
  \label{lem:error-control}
  Let $\pi_n$ be the sequence generated by~\eqref{eq:PG} with
  $\pi_1 \in \calU$. Then:
  \begin{enumerate}[label=\textup{(\alph*)}]
  \item $\calE_{n+1} \subseteq \calE_n$ and
    $H_{n+1} \subseteq H_n$ for all $n \ge 1$.
  \item $H_{n-1} \subseteq \calE_n$ for all $n \ge 1$.
  \item Consider the ``bad realisation'' event
    \begin{equation}
      \label{eq:bad-event}
      \tilde H_n \coloneqq H_{n-1} \setminus H_n
      = \{W_k \le a \text{ for } k = 1, \ldots, n-1
      \text{ and } W_n > a\},
    \end{equation}
    and set
    $\tilde W_n \coloneqq W_n \ind_{H_{n-1}}$. Then
    \begin{equation}
      \label{eq:Wtilde-recursion}
      \E[\tilde W_n] \le \E[\tilde W_{n-1}]
      + \gamma_n \norm{\Pi} B_n
      + \gamma_n^2 \norm{\Pi}^2 \sigma_n^2
      + \tfrac{3}{2}\gamma_n^2(G^2 + B_n^2 + \sigma_n^2)
      - a\, \Prob(\tilde H_{n-1}).
    \end{equation}
  \end{enumerate}
\end{lemma}

\begin{proof}
  Part~(a) is immediate from the definitions: both
  $\calE_n$ and $H_n$ are defined as intersections of events up to
  time~$n$, so adding one more constraint can only shrink them.

  For part~(b), assume $H_{n-1}$ holds, so $W_k \le a$ for all
  $k = 1, \ldots, n-1$. By telescoping the energy
  inequality~\eqref{eq:energy-ineq} and using
  $\ip{v(\pi_k)}{\pi_k - \pi^*} \le 0$ on $\calE_k$ (which holds when
  $\pi_k \in \calB$), we obtain, for any $k \le n$,
  \begin{equation}
    D_{k+1} \le D_1 + M_k + S_k
    \le a + \sqrt{W_k}
    \le a + \sqrt{a},
    \label{eq:D-bound-induction}
  \end{equation}
  where we used $D_1 \le a$ and
  $W_k = M_k^2 + S_k \ge M_k^2 \ge 0$, so
  $|M_k| + S_k \le |M_k| + M_k^2 + S_k - M_k^2 + S_k
  \le \sqrt{W_k} + W_k$; in fact, the key step is the bound from the
  energy inequality. Specifically, telescoping~\eqref{eq:energy-ineq}
  over $k = 1, \ldots, n$ and dropping the negative drift term gives
  \begin{equation}
    D_{n+1} \le D_1 + M_n + S_n.
    \label{eq:D-telescope}
  \end{equation}
  On $H_n$, we have $W_n = M_n^2 + S_n \le a$. Since $M_n^2 \ge 0$
  and $S_n \ge 0$, this implies $|M_n| \le \sqrt{a}$ and $S_n \le a$.
  Therefore
  $D_{n+1} \le a + \sqrt{a} + a = 2a + \sqrt{a} < \varrho^2/2$, so
  $\pi_{n+1} \in \calB$, confirming $\calE_{n+1}$. By induction,
  $H_{n-1} \subseteq \calE_n$.

  Part~(c) requires a detailed computation.
  We write $W_n - W_{n-1}$ explicitly using~\eqref{eq:W-def}:
  \begin{align}
    W_n &= (M_{n-1} + \gamma_n \xi_n)^2 + S_{n-1}
    + \gamma_n \chi_n + \gamma_n^2 \psi_n^2 \notag\\
    &= W_{n-1} + 2M_{n-1}\gamma_n \xi_n + \gamma_n^2 \xi_n^2
    + \gamma_n\chi_n + \gamma_n^2 \psi_n^2.
    \label{eq:W-increment}
  \end{align}
  Taking conditional expectations given $\calF_n$ (and noting that
  $M_{n-1}$ and $\ind_{H_{n-1}}$ are $\calF_n$-measurable, since they
  depend only on $\hat v_1, \ldots, \hat v_{n-1}$ and $\pi_1$):
  \begin{align}
    \E[W_n \mid \calF_n]
    &= W_{n-1} + 2M_{n-1}\gamma_n
    \underbrace{\E[\xi_n \mid \calF_n]}_{= 0}
    + \gamma_n^2 \E[\xi_n^2 \mid \calF_n]
    + \gamma_n\chi_n + \gamma_n^2 \E[\psi_n^2 \mid \calF_n]
    \notag \\
    &\ge W_{n-1},
    \label{eq:W-submart-check}
  \end{align}
  confirming that $W_n$ is a submartingale.

  Now define $\tilde W_n = W_n \ind_{H_{n-1}}$ and decompose:
  \begin{align}
    \tilde W_n &= W_n \ind_{H_{n-1}} \notag\\
    &= W_{n-1} \ind_{H_{n-1}} + (W_n - W_{n-1})\ind_{H_{n-1}} \notag\\
    &= W_{n-1} \ind_{H_{n-2}} - W_{n-1}\ind_{\tilde H_{n-1}}
      + (W_n - W_{n-1})\ind_{H_{n-1}},
    \label{eq:Wtilde-decomp}
  \end{align}
  where we used $H_{n-1} = H_{n-2} \setminus \tilde H_{n-1}$, so
  $\ind_{H_{n-1}} = \ind_{H_{n-2}} - \ind_{\tilde H_{n-1}}$. Thus
  \begin{equation}
    \tilde W_n = \tilde W_{n-1}
    + (W_n - W_{n-1})\ind_{H_{n-1}}
    - W_{n-1}\ind_{\tilde H_{n-1}}.
    \label{eq:Wtilde-recurrence}
  \end{equation}
  Taking expectations, we bound each piece:

  \emph{The noise term.}  Since $H_{n-1}$ and $M_{n-1}$ are both
  $\calF_n$-measurable, the tower property gives
  \begin{equation}
    \E[M_{n-1}\gamma_n\xi_n\, \ind_{H_{n-1}}]
    = \E\bigl[M_{n-1}\ind_{H_{n-1}}\,
      \gamma_n \E[\xi_n \mid \calF_n]\bigr] = 0.
    \label{eq:noise-vanishes}
  \end{equation}

  \emph{The squared noise.}  Using
  $\xi_n^2 = |\ip{U_n}{\pi_n-\pi^*}|^2
  \le \norm{U_n}^2 \norm{\pi_n-\pi^*}^2 \le \norm{\Pi}^2 \norm{U_n}^2$
  and $\E[\norm{U_n}^2 \mid \calF_n] \le \sigma_n^2$:
  \begin{equation}
    \E[\xi_n^2 \ind_{H_{n-1}}]
    \le \E\bigl[\ind_{H_{n-1}} \norm{\Pi}^2
      \E[\norm{U_n}^2 \mid \calF_n]\bigr]
    \le \norm{\Pi}^2 \sigma_n^2.
    \label{eq:xi-sq-bound}
  \end{equation}

  \emph{The variance term.}  Since
  $\norm{\hat v_n} \le \norm{v(\pi_n)} + \norm{U_n} + \norm{b_n}$ and
  $\norm{v(\pi_n)} \le G$, we have
  $\E[\psi_n^2 \ind_{H_{n-1}}]
  = \E[\tfrac{1}{2}\norm{\hat v_n}^2 \ind_{H_{n-1}}]
  \le \tfrac{3}{2}(G^2 + B_n^2 + \sigma_n^2)$.

  \emph{The bias term.}
  $\E[\chi_n\ind_{H_{n-1}}] \le \norm{\Pi} B_n$.

  \emph{The penalty.}  On $\tilde H_{n-1}$, we have $W_{n-1} > a$
  (since $\tilde H_{n-1} = H_{n-2} \setminus H_{n-1}$ means $W_{n-1}>a$
  while $W_k \le a$ for $k < n-1$), so
  $\E[W_{n-1}\ind_{\tilde H_{n-1}}] \ge a\, \Prob(\tilde H_{n-1})$.

  Combining these estimates with~\eqref{eq:Wtilde-recurrence} yields~\eqref{eq:Wtilde-recursion}.
\end{proof}

We now establish the main stability result.

\begin{proposition}[Stochastic stability;
  {\cite[Proposition~A.1]{giannou2022}}]
  \label{prop:stability}
  Fix $\delta > 0$ and let $\pi_n$ be generated by~\eqref{eq:PG} with
  step-size and gradient estimates as in Theorem~\ref{thm:main}. Then
  \begin{equation}
    \label{eq:stability}
    \Prob(H_n \mid \pi_1 \in \calU) \ge 1 - \delta
    \quad \text{for all } n = 1, 2, \ldots,
  \end{equation}
  provided $\gamma$ is sufficiently small (or $m$ sufficiently large)
  relative to $\delta$.
\end{proposition}

\begin{proof}
  We bound $\Prob(\tilde H_n)$ using the ``bad realisation''
  event~\eqref{eq:bad-event}. If $\pi_1 \in \calU$, then on
  $\tilde H_n$ we have $W_n > a$ and $W_k \le a$ for all $k < n$. Since
  $W_n \ge 0$, we have $\ind_{\tilde H_n} \le W_n \ind_{H_{n-1}} / a$,
  so
  \begin{equation}
    \Prob(\tilde H_n) = \E[\ind_{\tilde H_n}]
    \le \E[W_n \ind_{H_{n-1}}] / a
    = \E[\tilde W_n] / a.
    \label{eq:Htilde-Markov}
  \end{equation}

  Telescoping~\eqref{eq:Wtilde-recursion} from $1$ to $n$ gives
  \begin{equation}
    \E[\tilde W_n] \le \E[\tilde W_0]
    + \norm{\Pi} \sum_{k=1}^n \gamma_k B_k
    + \sum_{k=1}^n \gamma_k^2 \varrho_k^2
    - a \sum_{k=1}^n \Prob(\tilde H_{k-1}),
    \label{eq:Wtilde-telescope}
  \end{equation}
  where
  $\varrho_k^2 \coloneqq \norm{\Pi}^2 \sigma_k^2
  + \tfrac{3}{2}(G^2 + B_k^2 + \sigma_k^2)$ and
  $\tilde W_0 = W_0 = 0$ by convention. Therefore
  \begin{equation}
    \sum_{k=1}^n \Prob(\tilde H_k) \le \frac{1}{a}
    \sum_{k=1}^n \bigl[\gamma_k B_k \norm{\Pi}
      + \gamma_k^2 \varrho_k^2\bigr]
    \eqqcolon \frac{C}{a},
    \label{eq:sum-Htilde}
  \end{equation}
  where $C \equiv C(\gamma, m) > 0$ satisfies
  $\lim_{\gamma \to 0^+} C(\gamma, m) = 0$ and
  $\lim_{m \to \infty} C(\gamma, m) = 0$, by the summability conditions
  on $\gamma_n$, $B_n$, $\sigma_n$ (since $p + \ell_b > 1$ and
  $2p - 2\ell_\sigma > 1$).

  By choosing $\gamma$ small enough (or $m$ large enough), we can ensure
  $C/a < \delta$. Since the events $\tilde H_k$ are disjoint (each
  $\tilde H_k$ specifies the \emph{first} time $W$ exceeds $a$),
  \begin{equation}
    \Prob\Bigl(\bigcup_{k=1}^n \tilde H_k\Bigr)
    = \sum_{k=1}^n \Prob(\tilde H_k) \le \delta.
    \label{eq:union-Htilde}
  \end{equation}
  Since
  $H_n = \bigl(\bigcup_{k=1}^n \tilde H_k\bigr)^c \cap \{W_0 \le a\}$,
  we conclude
  $\Prob(H_n \mid \pi_1 \in \calU) \ge 1 - \delta$.
\end{proof}

\section{Extraction of a convergent subsequence}
\label{sec:conv-subseq}

\begin{proposition}[Convergent subsequence;
  {\cite[Proposition~A.2]{giannou2022}}]
  \label{prop:conv-subseq}
  Let $\pi^*$ be a stable Nash policy. Under the assumptions of
  Theorem~\ref{thm:main}, the sequence $\pi_n$ admits a subsequence
  $\pi_{n_k}$ converging to $\pi^*$ with probability~$1$ on the event
  $\calE \coloneqq \bigcap_n \calE_n = \{\pi_n \in \calB
  \text{ for all } n\}$.
\end{proposition}

\begin{proof}
  Define $\calQ \coloneqq \{\pi_n \in \calB \text{ for all } n\}
  \cap \{\liminf_n \norm{\pi_n - \pi^*} > 0\}$ --- the event that
  $\pi_n$ stays in $\calB$ but does \emph{not} admit a subsequence
  converging to $\pi^*$. We show $\Prob(\calQ) = 0$.

  Assume for contradiction that $\Prob(\calQ) > 0$. On $\calQ$, the
  stability of $\pi^*$ gives a (random but positive) constant $c > 0$
  such that
  $\ip{v(\pi_n)}{\pi_n - \pi^*} \le -c$ for all $n$. Summing the
  energy inequality~\eqref{eq:energy-ineq} from $k = 1$ to $n$ and
  using $D_n \ge 0$:
  \begin{equation}
    0 \le D_{n+1} \le D_1 - c\,\tau_n + M_n + S_n,
    \label{eq:D-contra-bound}
  \end{equation}
  where $\tau_n = \sum_{k=1}^n \gamma_k$. Since $p \le 1$, we have
  $\tau_n \to \infty$. Dividing by $\tau_n$:
  \begin{equation}
    0 \le \frac{D_1}{\tau_n} - c + \frac{M_n}{\tau_n}
    + \frac{S_n}{\tau_n}.
    \label{eq:contra-divided}
  \end{equation}

  \emph{$M_n/\tau_n \to 0$ a.s.}\; The martingale differences
  $d_n = \gamma_n \xi_n$ satisfy
  \begin{equation}
    \sum_{n=1}^\infty \frac{\E[d_n^2]}{\tau_n^2}
    \le \sum_{n=1}^\infty \frac{\gamma_n^2 \norm{\Pi}^2 \sigma_n^2}
    {\tau_n^2}.
    \label{eq:SLLN-check}
  \end{equation}
  Since $\gamma_n^2 \sigma_n^2 = \mathcal{O}(\gamma_n^2 \cdot n^{2\ell_\sigma})$
  and $p - \ell_\sigma > 1/2$ implies
  $\sum \gamma_n^2 \sigma_n^2 < \infty$ (as $2p - 2\ell_\sigma > 1$),
  and $\tau_n \to \infty$, the strong law for martingale differences
  (Theorem~\ref{thm:SLLN-mart}) gives $M_n/\tau_n \to 0$ a.s.

  \emph{$S_n/\tau_n \to 0$ a.s.}\; By the bounds on $B_n$ and
  $\sigma_n$,
  \begin{equation}
    \E[S_n] = \sum_{k=1}^n \bigl[\gamma_k \chi_k
      + \gamma_k^2 \E[\psi_k^2]\bigr]
    \le \norm{\Pi}\sum_{k=1}^n \gamma_k B_k
    + \tfrac{3}{2}\sum_{k=1}^n \gamma_k^2(G^2 + B_k^2 + \sigma_k^2).
    \label{eq:ES-bound}
  \end{equation}
  Since $p + \ell_b > 1$, the first sum converges, and since
  $2p - 2\ell_\sigma > 1$, the second converges as well. Hence $S_n$
  is bounded in $L^1$, and by Doob's submartingale convergence theorem
  (Theorem~\ref{thm:doob-conv}), $S_n \to S_\infty$ a.s.\ for some
  finite $S_\infty$. In particular, $S_n / \tau_n \to 0$.

  Letting $n \to \infty$ in~\eqref{eq:contra-divided},
  $0 \le -c + 0 + 0 = -c < 0$, a contradiction. Since $D \ge 0$
  and $\Prob(\calQ) > 0$ was assumed, we conclude $\Prob(\calQ) = 0$.
\end{proof}

\section{Convergence of the energy function}
\label{sec:energy-conv}

\begin{proposition}[Energy convergence;
  {\cite[Proposition~A.3]{giannou2022}}]
  \label{prop:energy-conv}
  Under the assumptions of Proposition~\ref{prop:stability}, there
  exists a finite random variable $D_\infty$ such that
  \begin{equation}
    \label{eq:D-convergence}
    \Prob\bigl(D_n \to D_\infty \text{ as } n \to \infty
    \mid \pi_n \in \calB \text{ for all } n\bigr) = 1.
  \end{equation}
\end{proposition}

\begin{proof}
  Let $\calE_n = \{\pi_k \in \calB \text{ for all } k = 1, \ldots, n\}$
  as before, and define
  $\tilde D_n \coloneqq \ind_{\calE_n}\, D_n$. Since
  $\calE_{n+1} \subseteq \calE_n$ (Lemma~\ref{lem:error-control}(a)), we
  have $\ind_{\calE_{n+1}} \le \ind_{\calE_n}$, and so
  \begin{align}
    \tilde D_{n+1}
    &= \ind_{\calE_{n+1}}\, D_{n+1}
    \le \ind_{\calE_n}\, D_{n+1}.
    \label{eq:Dtilde-bound}
  \end{align}
  On $\calE_n$, the energy inequality~\eqref{eq:energy-ineq} combined
  with the stability condition
  $\ip{v(\pi_k)}{\pi_k - \pi^*} \le 0$ gives
  \begin{equation}
    \ind_{\calE_n}\, D_{n+1}
    \le \ind_{\calE_n}\, D_n
    + \gamma_n\, \ind_{\calE_n}\, \xi_n
    + \ind_{\calE_n}\bigl(\gamma_n \chi_n + \gamma_n^2 \psi_n^2\bigr).
    \label{eq:Dtilde-ineq}
  \end{equation}
  Thus
  $\tilde D_{n+1} \le \tilde D_n
  + \gamma_n\, \ind_{\calE_n}\, \xi_n
  + \gamma_n \chi_n + \gamma_n^2 \psi_n^2$.
  Taking conditional expectations given~$\calF_n$:
  \begin{align}
    \E[\tilde D_{n+1} \mid \calF_n]
    &\le \tilde D_n + \gamma_n \ind_{\calE_n}
      \underbrace{\E[\xi_n \mid \calF_n]}_{= 0}
      + \gamma_n \chi_n
      + \gamma_n^2 \E[\psi_n^2 \mid \calF_n]
    \notag\\
    &\le \tilde D_n + \gamma_n\, \norm{\Pi}\, B_n
      + \tfrac{3}{2}\gamma_n^2 (G^2 + B_n^2 + \sigma_n^2).
    \label{eq:Dtilde-cond-exp}
  \end{align}
  Denoting the nonnegative remainder by
  $a_n \coloneqq \gamma_n\, \norm{\Pi}\, B_n
  + \tfrac{3}{2}\gamma_n^2 (G^2 + B_n^2 + \sigma_n^2)$, we have
  $\sum_n a_n < \infty$ by the step-size and bias/variance assumptions.
  The Robbins--Siegmund lemma (Theorem~\ref{thm:robbins-siegmund}),
  applied with $X_n = \tilde D_n$, $b_n = 0$, and $a_n$ as above,
  yields that $\tilde D_n$ converges almost surely to a finite random
  variable $\tilde D_\infty$.

  On the event $\calE = \bigcap_n \calE_n$, we have
  $\ind_{\calE_n} = 1$ for all $n$, so $\tilde D_n = D_n$ and hence
  $D_n \to \tilde D_\infty$ a.s.\ on $\calE$. Setting
  $D_\infty = \tilde D_\infty$ on $\calE$ gives the result.
\end{proof}

\section{Proof of the main theorem}
\label{sec:final-proof}

We are now in a position to prove Theorem~\ref{thm:main}.

\begin{proof}[Proof of Theorem~\ref{thm:main}]
  Let $\calE = \bigcap_n \calE_n = \{\pi_n \in \calB \text{ for all }
  n\}$. By Proposition~\ref{prop:stability} (via
  Lemma~\ref{lem:error-control}(b), which gives
  $H_n \subseteq \calE_{n+1}$), we have
  \begin{equation}
    \Prob(\calE \mid \pi_1 \in \calU) \ge 1 - \delta,
    \label{eq:E-prob}
  \end{equation}
  noting that the neighbourhood $\calU$ is independent of the
  confidence level $\delta$.

  By Propositions~\ref{prop:conv-subseq} and~\ref{prop:energy-conv},
  the following hold with probability $1$ on~$\calE$:
  \begin{enumerate}[label=(\roman*)]
  \item $\liminf_n \norm{\pi_n - \pi^*} = 0$ (there is a subsequence
    converging to $\pi^*$).
  \item $D_n$ converges to a finite limit $D_\infty$.
  \end{enumerate}
  Since $D_n = \tfrac{1}{2}\norm{\pi_n - \pi^*}^2$, statement~(i)
  implies that some subsequence $D_{n_k} \to 0$.
  Statement~(ii) says $D_n \to D_\infty$. A convergent sequence with a
  subsequence converging to zero must itself converge to zero: hence
  $D_\infty = 0$ and $\pi_n \to \pi^*$.

  Therefore
  \begin{align}
    \Prob(\pi_n \to \pi^* \mid \pi_1 \in \calU)
    &\ge \Prob(\calE \cap \{\pi_n \to \pi^*\} \mid \pi_1 \in \calU)
      \notag\\
    &= \Prob(\pi_n \to \pi^* \mid \pi_1 \in \calU,\, \calE)\;
       \Prob(\calE \mid \pi_1 \in \calU)
    \notag\\
    &\ge 1 \cdot (1 - \delta) = 1 - \delta,
    \label{eq:final-bound}
  \end{align}
  completing the proof.
\end{proof}


\section{Application to the REINFORCE algorithm}
\label{sec:corollary}

We conclude this chapter by recording the consequence of
Theorem~\ref{thm:main} for the three gradient information models
described in Section~\ref{sec:PG}.

\begin{corollary}[{\cite[Corollary~1]{giannou2022}}]
  \label{cor:models}
  Suppose that~\eqref{eq:PG} is run with step-size
  $\gamma_n = \gamma/(n+m)^p$, $p > 1/2$, and, if applicable, an
  exploration parameter $\varepsilon_n = \varepsilon/(n+m)^r$ with
  $1-p < r < 2p - 1$. Then:
  \begin{itemize}
  \item \emph{Full gradients} ($\hat v_n = v(\pi_n)$, so $B_n = 0$,
    $\sigma_n = 0$): the conclusions of Theorem~\ref{thm:main} hold for
    any $p > 1/2$.
  \item \emph{Stochastic gradients} ($B_n = 0$,
    $\sigma_n = \mathcal{O}(1)$): same conclusion, with
    $p > 1/2$.
  \item \emph{$\varepsilon$-Greedy REINFORCE}
    ($B_n = \mathcal{O}(\varepsilon_n)$,
    $\sigma_n = \mathcal{O}(1/\sqrt{\varepsilon_n})$): the conclusion
    holds for $p > 2/3$.
  \end{itemize}
\end{corollary}

\begin{proof}
  For full and stochastic gradients, $\ell_b = \infty$ and
  $\ell_\sigma = 0$, so the conditions $p + \ell_b > 1$ and
  $p - \ell_\sigma > 1/2$ reduce to $p > 1/2$.

  For the $\varepsilon$-Greedy variant, $B_n = \mathcal{O}(\varepsilon_n)$
  with $\varepsilon_n = \varepsilon/(n+m)^r$, giving $\ell_b = r$.
  The variance satisfies
  $\sigma_n = \mathcal{O}(1/\sqrt{\varepsilon_n})
  = \mathcal{O}(n^{r/2})$, giving $\ell_\sigma = r/2$. The conditions
  become $p + r > 1$ and $p - r/2 > 1/2$, i.e.,
  $r > 1 - p$ and $r < 2p - 1$. For such an $r$ to exist, we need
  $1 - p < 2p - 1$, i.e., $p > 2/3$.
\end{proof}


%% ================================================================
%%  CHAPTER 4: GLOBAL CONVERGENCE
%% ================================================================
\chapter{Global Convergence via Random Restarts}
\label{ch:global}

Theorem~\ref{thm:main} guarantees convergence to a stable Nash policy
provided the initial policy lies in a neighbourhood $\calU$ of the
equilibrium. In practice, agents do not know where Nash policies are
located, so the dependence on initialisation is a serious limitation.
This chapter removes it: we prove that a simple restart strategy
converges to a Nash equilibrium \emph{almost surely}, regardless of
how the game is structured, provided at least one stable Nash policy
exists.

\section{The restart policy gradient algorithm}
\label{sec:restart-alg}

The idea is to run independent PG trajectories from fresh random
initialisations. If any trajectory starts inside the basin of
attraction, Theorem~\ref{thm:main} ensures it converges.

\begin{definition}[Restart-PG]
  \label{def:restart-PG}
  Fix step-size parameters $\gamma, m, p$ satisfying the conditions of
  Theorem~\ref{thm:main}. A \emph{restart schedule} is a sequence of
  independent runs $(R_k)_{k \ge 1}$, where each run~$R_k$ consists of:
  \begin{enumerate}
  \item An initialisation $\pi_1^{(k)}$ drawn independently from a
    distribution $\nu$ on $\Pi$.
  \item A PG trajectory $(\pi_n^{(k)})_{n \ge 1}$ generated
    by~\eqref{eq:PG} from $\pi_1^{(k)}$, using independent episode
    trajectories.
  \end{enumerate}
  We say epoch $k$ is \emph{successful} if $\pi_n^{(k)}$ converges to
  some stable Nash policy as $n \to \infty$.
\end{definition}

\section{Product probability space}
\label{sec:product-space}

The measure-theoretic foundation for the restart argument is a product
construction. Each epoch $k$ is supported on a probability space
$(\Omega_k, \calF_k, \Prob_k)$ as constructed in
Section~\ref{sec:prob-framework}, carrying the initialisation
$\pi_1^{(k)}$ and all episode trajectories
$\tau_1^{(k)}, \tau_2^{(k)}, \ldots$ of that epoch. The full restart
process lives on the product space
\begin{equation}
  \label{eq:product-space}
  (\Omega, \calF, \Prob) = \prod_{k=1}^\infty
  (\Omega_k, \calF_k, \Prob_k),
\end{equation}
constructed via the Kolmogorov extension theorem~\cite{durrett2019}.
Since each factor is independent, the success events
\begin{equation}
  \label{eq:success-event}
  A_k \coloneqq \bigl\{\omega \in \Omega :
  \pi_n^{(k)}(\omega) \to \pi^*
  \text{ for some stable Nash } \pi^*\bigr\}
\end{equation}
are mutually independent: $A_k$ depends only on the $k$-th coordinate
of the product space.

\section{Almost-sure convergence}
\label{sec:as-convergence}

We make the following assumptions.

\begin{enumerate}[label=(A\arabic*)]
\item \label{ass:stable-exists} The game $\calG$ possesses at least one
  strict Nash policy (Definition~\ref{def:refinements}).
  Since every strict Nash policy is stable, the local convergence
  guarantee of Theorem~\ref{thm:main} applies at each such
  equilibrium.
\item \label{ass:full-support} The initialisation distribution $\nu$ is
  absolutely continuous with respect to the Lebesgue measure
  $\lambda_\Pi$ on $\Pi$, with Radon--Nikodym density bounded below:
  $d\nu / d\lambda_\Pi \ge c$ for some $c > 0$. (The uniform
  distribution on $\Pi$ satisfies this with
  $c = 1/\lambda_\Pi(\Pi)$.)
\end{enumerate}

Let $\pi^{*,1}, \ldots, \pi^{*,J}$ denote the stable Nash policies of
$\calG$, with corresponding neighbourhoods
$\calU_1, \ldots, \calU_J \subseteq \Pi$ from Theorem~\ref{thm:main}.
Define
\begin{equation}
  \label{eq:basin-prob}
  p \coloneqq \nu\Bigl(\bigcup_{j=1}^J \calU_j\Bigr) > 0,
\end{equation}
which is strictly positive because each $\calU_j$ is a ball of positive
radius in $\Pi$ and $\nu$ has full support.

\begin{theorem}[Global convergence via restarts]
  \label{thm:global}
  Under assumptions~\ref{ass:stable-exists}
  and~\ref{ass:full-support}, choose step-size parameters $\gamma, m$
  so that Theorem~\ref{thm:main} holds with some confidence level
  $\delta \in (0,1)$. Then:
  \begin{enumerate}[label=\textup{(\alph*)}]
  \item The restart schedule converges to a Nash equilibrium almost
    surely:
    \begin{equation}
      \label{eq:global-as}
      \Prob\bigl(\exists\, k \ge 1 : A_k \text{ occurs}\bigr) = 1.
    \end{equation}
  \item The index $K^*$ of the first successful epoch satisfies
    \begin{equation}
      \label{eq:geom-bound}
      \Prob(K^* > k) \le \bigl(1 - p(1-\delta)\bigr)^k
      \quad \text{for all } k \ge 1,
    \end{equation}
    and in particular $\E[K^*] \le 1/\bigl(p(1-\delta)\bigr)$.
  \end{enumerate}
\end{theorem}

\begin{proof}
  By Theorem~\ref{thm:main}, for each epoch $k$,
  \begin{equation}
    \Prob(A_k) \ge \Prob\bigl(\pi_1^{(k)} \in \textstyle\bigcup_j
    \calU_j\bigr) \cdot \min_j \Prob\bigl(\pi_n^{(k)} \to \pi^{*,j}
    \bigm| \pi_1^{(k)} \in \calU_j\bigr)
    \ge p(1-\delta).
    \label{eq:epoch-success}
  \end{equation}
  Since the events $A_k$ are independent
  (Section~\ref{sec:product-space}), each with
  $\Prob(A_k) \ge p(1-\delta) > 0$:
  \begin{equation}
    \Prob\bigl(\textstyle\bigcap_{k=1}^K A_k^c\bigr)
    = \prod_{k=1}^K \Prob(A_k^c)
    \le \bigl(1 - p(1-\delta)\bigr)^K.
    \label{eq:no-success}
  \end{equation}
  Letting $K \to \infty$ and using
  $1 - p(1-\delta) < 1$:
  \begin{equation}
    \Prob\bigl(\textstyle\bigcap_{k=1}^\infty A_k^c\bigr)
    = \lim_{K \to \infty} \bigl(1 - p(1-\delta)\bigr)^K = 0.
    \label{eq:as-convergence}
  \end{equation}
  This proves~\eqref{eq:global-as}. The bound~\eqref{eq:geom-bound}
  is~\eqref{eq:no-success} with $K^* > k$ iff no epoch among
  $1, \ldots, k$ succeeds.
\end{proof}

\section{Restart complexity}
\label{sec:restart-complexity}

The expected number of restarts $\E[K^*] = 1/(p(1-\delta))$ is
controlled by two quantities: the confidence level $\delta$ (which can
be made arbitrarily small) and the basin probability~$p$.

\paragraph{Basin volume.}
From the proof of Theorem~\ref{thm:main}
(Section~\ref{sec:stability}), the neighbourhood $\calU_j$ is a ball of
radius $r_j = \sqrt{2a_j}$ centred at $\pi^{*,j}$, where $a_j$
satisfies $2a_j + \sqrt{a_j} < \varrho_j$ and $\varrho_j$ is the
stability radius (the largest $\varrho$ such that
$\ip{v(\pi)}{\pi - \pi^{*,j}} < 0$ for all
$\pi \in \calB(\pi^{*,j}, \varrho) \setminus \{\pi^{*,j}\}$). For an
SOS policy with parameter $\mu$ and gradient Lipschitz constant $L$,
the stability radius satisfies $\varrho \ge \mu / L$, giving
$a = \Theta(\mu^2/L^2)$ and $r = \Theta(\mu/L)$.

When $\nu$ is uniform on $\Pi$, the basin probability is
\begin{equation}
  \label{eq:basin-volume}
  p \ge \frac{\lambda_\Pi\bigl(\calU_j \cap \Pi\bigr)}
  {\lambda_\Pi(\Pi)},
\end{equation}
where the numerator is the volume of a ball of radius
$\Theta(\mu/L)$ intersected with the policy simplex, in dimension
$d = |\calS| \sum_i (|\calA_i| - 1)$. This gives
\begin{equation}
  \label{eq:restart-bound}
  \E[K^*] = \mathcal{O}\!\left(
    \frac{\lambda_\Pi(\Pi)}{\omega_d \cdot (\mu/L)^d \cdot (1-\delta)}
  \right),
\end{equation}
where $\omega_d$ is the volume of the unit ball in $\R^d$.

\begin{remark}[Effective dimension and the blessing of dimensionality]
  \label{rem:dimension}
  The worst-case bound~\eqref{eq:restart-bound} scales exponentially in
  the ambient dimension $d$ of the policy space, which might suggest
  that random restarts are impractical for large games. However, this
  bound is pessimistic: it treats the basin $\calU_j$ as a ball of
  fixed radius in all $d$ directions. In practice, the gradient field
  $v(\pi)$ near a Nash policy typically has low effective rank---the
  Jacobian $\operatorname{Jac}_v(\pi^*)$ has a decaying spectrum,
  concentrating the dynamics on a low-dimensional subspace. By
  concentration of measure in high dimensions, the basin of attraction
  may extend much further in the ``irrelevant'' directions (those
  corresponding to small eigenvalues) than in the ``relevant'' ones.
  The effective restart complexity then scales with the
  \emph{effective dimension}
  $d_{\mathrm{eff}} = \operatorname{tr}(\Sigma)/\lambda_{\max}(\Sigma)$
  rather than $d$, where $\Sigma$ is the covariance of the gradient
  noise. When the optimal policy is $k$-sparse ($k \ll d$), one has
  $d_{\mathrm{eff}} \approx k$ and the restart cost scales as
  $\mathcal{O}(k \log d)$---sub-linear in the action space size.
  Structured exploration methods (homotopy continuation, Bayesian
  optimisation) can reduce restart counts further.
\end{remark}

\begin{remark}[Computational hardness]
  \label{rem:ppad}
  Foster et~al.~\cite{foster2023} proved that independent no-regret
  learning in general-sum Markov games requires exponentially many
  episodes in the worst case (PPAD-hardness). Restart-PG does not
  contradict this: its efficiency depends on the
  \emph{instance-dependent} basin probability $p$, not worst-case
  structure. The framework shifts complexity from worst-case to
  instance-dependent---the appropriate regime when equilibria are
  sparse (Remark~\ref{rem:dimension}).
\end{remark}

\begin{remark}[Distinction from stochastic resetting]
  \label{rem:resetting}
  \emph{Stochastic resetting}---returning an agent to a reference state
  mid-episode---is a micro-level, intra-episodic temporal truncation
  tool. Restart-PG is a macro-level, inter-episodic spatial exploration
  mechanism that re-initialises the entire policy vector. The two are
  complementary.
\end{remark}

\section{Equilibrium selection}
\label{sec:eq-selection}

When $\calG$ has multiple stable Nash policies, different successful
epochs may converge to different equilibria. The restart framework
naturally enables \emph{equilibrium selection}: after $K$ epochs, we
compare the discovered equilibria and select the one maximising a
welfare criterion.

\begin{definition}[Social welfare]
  For a Nash policy $\pi^*$, define the social welfare
  $\mathrm{SW}(\pi^*) = \sum_{i \in \calN} V_{i,\rho}(\pi^*)$.
\end{definition}

Suppose $K$ restarts have discovered distinct Nash policies
$\hat\pi^{(1)}, \ldots, \hat\pi^{(m)}$ (with $m \le K$). We select
$\hat\pi^{(\text{best})} = \argmax_{j \le m} \mathrm{SW}(\hat\pi^{(j)})$.

\begin{proposition}[Discovery probability]
  \label{prop:discovery}
  If Nash policy $\pi^{*,j}$ has basin probability $p_j > 0$, then after
  $K$ restarts,
  \begin{equation}
    \Prob\bigl(\pi^{*,j} \text{ is discovered}\bigr)
    \ge 1 - \bigl(1 - p_j(1-\delta)\bigr)^K.
  \end{equation}
  In particular, the welfare-maximising Nash policy is discovered with
  probability at least $1 - (1 - p_{\max}(1-\delta))^K$, where
  $p_{\max}$ is the basin probability of the best equilibrium.
\end{proposition}

\begin{proof}
  Each restart independently discovers $\pi^{*,j}$ with probability at
  least $p_j(1-\delta)$. The complementary probability over $K$
  independent trials gives the bound.
\end{proof}

This provides a principled stopping criterion: continue restarting until
the marginal probability of discovering a strictly better Nash policy
falls below a threshold. In practice, this can be estimated using
species-discovery methods (such as the Chao1 estimator) applied to the
set of discovered equilibria.


\section{Empirical validation}
\label{sec:simulations}

We validate Theorem~\ref{thm:global} on three two-player normal-form
games with multiple Nash equilibria: Coordination, Stag Hunt, and
Battle of the Sexes (payoff matrices are given in
Appendix~\ref{app:games}). In each experiment, agents use projected
policy gradient with step-size $\gamma_n = 0.3/(n+10)^{0.75}$ and
REINFORCE-style noise of varying magnitude. The simulator was coded
from scratch in Python (NumPy); each heatmap in
Figure~\ref{fig:basins} evaluates a $35 \times 35$ grid of
initialisations with $2{,}000$ gradient steps per point, and all
random experiments use seed $42$ for reproducibility. The first three
experiments specialise to normal-form games ($|\calS| = 1$) for
clarity; a full stochastic game validation ($|\calS| = 2$) follows.

\paragraph{Basins of attraction.}
Figure~\ref{fig:basins} visualises the basins of attraction under
deterministic PG. Each point $(\pi_1(a_1), \pi_2(a_1)) \in [0,1]^2$
represents an initialisation; the colour indicates which Nash
equilibrium the trajectory converges to. The basins tile nearly the
entire strategy space, separated by a thin boundary. This explains the
key empirical finding: the basin probability $p$ from
Theorem~\ref{thm:global} is close to~$1$ for these games, and a single
random restart almost always suffices for convergence.

\begin{figure}[ht]
  \centering
  \includegraphics[width=\textwidth]{basins_of_attraction.pdf}
  \caption{Basins of attraction under deterministic PG for three
    two-player games. Stars mark Nash equilibria. Red and blue
    regions converge to distinct Nash policies; grey indicates
    non-convergence. The basins cover nearly the entire strategy space,
    confirming that the basin probability $p \approx 1$ in
    Theorem~\ref{thm:global}, so $K^* = 1$ with overwhelming
    probability (Remark~\ref{rem:dimension}).}
  \label{fig:basins}
\end{figure}

\paragraph{Equilibrium selection.}
While convergence to \emph{some} Nash is easy (one restart suffices),
finding the \emph{best} Nash requires multiple restarts. Figure~\ref{fig:selection}
shows this for Stag Hunt, which has two Nash equilibria with social
welfare $\mathrm{SW} = 8$ (Stag, Stag) and $\mathrm{SW} = 4$ (Hare,
Hare). A single restart yields average welfare around $5$ (often
landing on the risk-dominant but welfare-inferior equilibrium). With
$15$ restarts, the best-of-$K$ welfare approaches the Pareto optimum.
This confirms Proposition~\ref{prop:discovery}: restarts serve
primarily as an equilibrium \emph{selection} mechanism, not merely a
convergence mechanism.

\begin{figure}[ht]
  \centering
  \includegraphics[width=0.75\textwidth]{equilibrium_selection.pdf}
  \caption{Stag Hunt: social welfare of the best discovered Nash
    policy as a function of the number of restarts $K$. The shaded
    region shows $\pm 1$ standard deviation over $100$ trials. Welfare
    approaches the Pareto-optimal equilibrium as $K$ increases,
    confirming Proposition~\ref{prop:discovery}: restarts serve as an
    equilibrium \emph{selection} mechanism.}
  \label{fig:selection}
\end{figure}

\paragraph{Dimension scaling.}
We also tested restart complexity on random $d \times d$ games with
diagonal-dominant payoffs (ensuring sparse Nash equilibria exist), for
$d = 2, 3, 4, 5, 7$. In every game tested, a single restart sufficed:
$K^* = 1$ regardless of dimension. This is consistent with the blessing
of dimensionality discussed in Remark~\ref{rem:dimension}: when Nash
policies are sparse, the basins of attraction extend broadly in the
irrelevant action dimensions, and the effective restart complexity
depends on the number of actions that matter, not the total action space
size.

\paragraph{Stochastic game validation ($|\calS| = 2$).}
To validate beyond the $|\calS| = 1$ setting, we test on a two-state
\emph{state-dependent coordination} game
(Appendix~\ref{app:games}): in state~$s_1$ the optimal play is
$a_1$; in state~$s_2$ it is $a_2$. The game has multiple strict Nash
policies: the state-contingent optimum with $\mathrm{SW} = 5.0$ and
two state-blind Nash with $\mathrm{SW} \approx 3.1$.
Figure~\ref{fig:stoch-convergence} shows convergence of episodic PG
with REINFORCE gradients from random initialisations: the agent learns
state-contingent play, with $\pi_1(a_1 \mid s_1) \to 1$ and
$\pi_1(a_1 \mid s_2) \to 0$.

\begin{figure}[ht]
  \centering
  \includegraphics[width=\textwidth]{simulations/figures/stochastic_pg/stochastic_convergence.pdf}
  \caption{State-Dependent Coordination ($|\calS| = 2$): convergence
    of episodic PG with REINFORCE gradients. Left: distance to
    nearest Nash from $8$ random initialisations. Right: policy
    probabilities $\pi_1(a_1 \mid s)$ for both states, showing the
    emergence of state-contingent play.}
  \label{fig:stoch-convergence}
\end{figure}

Figure~\ref{fig:stoch-comparison} shows that state structure matters:
state-contingent PG achieves $\mathrm{SW} = 5.0$, while
state-blind PG (the $|\calS| = 1$ reduction) is stuck at
$\mathrm{SW} \approx 3.1$ --- a $38\%$ welfare gap.

\begin{figure}[ht]
  \centering
  \includegraphics[width=\textwidth]{simulations/figures/stochastic_pg/stochastic_vs_normal.pdf}
  \caption{State-contingent vs.\ state-blind PG. Left: social welfare
    over iterations (mean $\pm$ s.d.\ over $20$ runs). Right: final
    welfare at convergence. The state-contingent agent achieves the
    theoretical optimum; the state-blind agent converges to a
    strictly inferior equilibrium.}
  \label{fig:stoch-comparison}
\end{figure}


%% ================================================================
%%  CHAPTER 5: CONCLUSION
%% ================================================================
\chapter{Concluding Remarks}
\label{ch:conclusion}

We have presented two results on the convergence of policy gradient
methods in multi-agent stochastic games.

The first (Chapter~\ref{ch:proof}) is a self-contained,
measure-theoretic exposition of Theorem~1 of Giannou
et~al.~\cite{giannou2022}: if all agents run policy gradient from a
neighbourhood of a stable Nash policy, the sequence of play converges to
the equilibrium with high probability. The proof rests on a Lyapunov
energy inequality, martingale stability estimates, and the
Robbins--Siegmund almost supermartingale convergence lemma.

The second (Chapter~\ref{ch:global}) is an original result: by running
independent PG trajectories from random initialisations, convergence to
a Nash equilibrium is achieved \emph{almost surely}. The expected number
of restarts is controlled by the basin volume of the equilibria, and the
framework naturally enables equilibrium selection among multiple Nash
policies.

Together, these results give a complete picture: local convergence
guarantees provide the engine, and random restarts provide the ignition.
The combination transforms a local result into a global one, at the cost
of a restart complexity that scales with the dimension of the policy
space.

\paragraph{Positioning within the 2025--2026 landscape.}
Table~\ref{tab:comparison} contrasts Restart-PG with recent
approaches. Alternating PG~\cite{kalogiannis2025} and Convex Markov
Game frameworks achieve global convergence through structural
innovations (Proximal Polyak--{\L}ojasiewicz, agent-wise gradient
domination) but converge to a single $\varepsilon$-approximate
equilibrium without selection capability. Restart-PG is complementary:
it circumvents non-convexity through repeated geometric sampling,
with total sample complexity
$\mathcal{O}(K^* \cdot T(\varepsilon))$. The overhead $K^*$ is the
mandatory cost of two features no other method provides: \emph{exact}
equilibrium fidelity and \emph{Pareto-optimal selection}.

\begin{table}[ht]
  \centering
  \small
  \caption{Comparison of global convergence methods for stochastic
    games. Restart-PG is the only method that natively supports
    equilibrium selection.}
  \label{tab:comparison}
  \begin{tabular}{lcccc}
    \toprule
    \textbf{Method} & \textbf{Mechanism} &
    \textbf{Equilibrium} & \textbf{Model-free?} &
    \textbf{Selection?} \\
    \midrule
    Restart-PG & Spatial exploration &
      Exact stable NE & Yes & \textbf{Yes} \\
    Alternating PG & Proximal-P{\L} + &
      $\varepsilon$-approx.\ & Yes & No \\
    & asym.\ timescales & (regularised) & & \\
    Convex MG & Agent-wise gradient &
      $\varepsilon$-approx.\ & No & No \\
    & domination & & (gen.\ model) & \\
    \bottomrule
  \end{tabular}
\end{table}

\paragraph{Limitations.}
Our experiments use normal-form games and a single $|\calS| = 2$
stochastic game. The stochastic game experiment confirms that the
theory applies beyond the $|\calS| = 1$ setting, but each restart
requires full episode rollouts, so the effective cost depends on the
horizon length and the dimensionality of the policy space.
The restart bound~\eqref{eq:restart-bound} suffers from the curse of
dimensionality; structured exploration
(homotopy continuation, Bayesian optimisation over the simplex) could
reduce this substantially.

\paragraph{Future work.}
Several directions remain open.
\begin{enumerate}[label=(\alph*)]
\item \emph{Function approximation.} Extending the restart framework
  beyond tabular policies to neural network parameterisations would
  broaden applicability, though the basin-of-attraction geometry
  becomes far more complex.
\item \emph{Optimistic and regularised PG variants.} Recent work on
  optimistic policy gradient in extensive-form
  games~\cite{kalogiannis2025} and softmax PG as a multidimensional
  Wiener process~\cite{leung2024} suggests that refined update rules
  may enlarge basins of attraction, reducing the restart count.
  Equilibrium collapse phenomena in single-controller
  games~\cite{anagnostides2024} indicate that stability is not
  universal, reinforcing the need for the strict-equilibrium assumption
  adopted here.
\item \emph{Non-stationary environments.} When the game's payoff
  structure changes over time, the restart mechanism could be adapted
  to a sliding-window or discounted scheme that tracks moving
  equilibria.
\item \emph{Asymmetric learning.} The case where only some agents use
  policy gradient while others use different rules (e.g., Q-learning)
  requires fundamentally different techniques, since the gradient
  dominance property relies on all agents updating via policy gradients
  simultaneously.
\item \emph{Convergence rates.} Combining the
  $\mathcal{O}(1/\sqrt{n})$ rate from Giannou's Theorem~2 with the
  geometric restart distribution would yield end-to-end sample
  complexity bounds.
\end{enumerate}

\bigskip

\noindent\textbf{Acknowledgements.} I am grateful to Prof.\ Galit
Ashkenazi-Golan for her supervision and guidance throughout this
project.

%% ================================================================
%%  APPENDIX: GAME DEFINITIONS
%% ================================================================
\appendix
\chapter{Experimental Game Definitions}
\label{app:games}

This appendix records the payoff matrices and transition kernels for
the games used in Section~\ref{sec:simulations}.

\section*{Normal-form games ($|\calS| = 1$)}

Each game below is a two-player, two-action game with actions
$\calA_1 = \calA_2 = \{a_1, a_2\}$. Entries are $(R_1, R_2)$
where $R_i$ is the payoff to player~$i$.

\paragraph{Coordination.}
Both players prefer to coordinate on the same action, with
$a_1$ Pareto-dominant.
\[
\begin{pmatrix} 2,\;2 & 0,\;0 \\ 0,\;0 & 1,\;1 \end{pmatrix}
\]
Nash equilibria: $(a_1, a_1)$ with $\mathrm{SW} = 4$;
$(a_2, a_2)$ with $\mathrm{SW} = 2$; and a mixed equilibrium at
$\pi_i(a_1) = 1/3$. Both pure equilibria are strict.

\paragraph{Stag Hunt.}
Models the tension between a payoff-dominant equilibrium
(Stag, Stag) and a risk-dominant one (Hare, Hare).
\[
\begin{pmatrix} 4,\;4 & 0,\;3 \\ 3,\;0 & 2,\;2 \end{pmatrix}
\]
Nash equilibria: $(a_1, a_1)$ with $\mathrm{SW} = 8$ (Stag);
$(a_2, a_2)$ with $\mathrm{SW} = 4$ (Hare); and a mixed equilibrium
at $\pi_i(a_1) = 2/3$. Both pure equilibria are strict.

\paragraph{Battle of the Sexes.}
An anti-coordination game with asymmetric payoffs.
\[
\begin{pmatrix} 3,\;2 & 0,\;0 \\ 0,\;0 & 2,\;3 \end{pmatrix}
\]
Nash equilibria: $(a_1, a_1)$ with $\mathrm{SW} = 5$;
$(a_2, a_2)$ with $\mathrm{SW} = 5$; and a mixed equilibrium at
$\pi_1(a_1) = 3/5$, $\pi_2(a_1) = 2/5$. Both pure equilibria
are strict.

\paragraph{Dimension-scaling games.}
For the scaling experiment (Section~\ref{sec:simulations},
Figure~\ref{fig:basins}), we generate random $d \times d$ games with
diagonal-dominant payoffs:
$R_i = 2\,I_d + 0.3\,Z_i$, where $Z_i$ has i.i.d.\
$\mathcal{N}(0,1)$ entries. The diagonal dominance ensures at least
one strict pure Nash equilibrium exists ($a_k = a_k$ for the action
$k$ maximising the diagonal), while the off-diagonal noise creates
a non-trivial landscape.

\section*{Stochastic game ($|\calS| = 2$): State-Dependent Coordination}

A two-player, two-action, two-state episodic game with termination
probability $\zeta = 0.4$ and uniform initial distribution
$\rho = (1/2, 1/2)$. In state~$s_1$ the preferred coordination is
on $a_1$; in state~$s_2$ the preferred coordination is on $a_2$.

\paragraph{Payoff tensors $R_i(s, a_1^{}, a_2^{})$} (identical for
both players):
\[
R(s_1) = \begin{pmatrix} 1.0 & -0.5 \\ -0.5 & 0.3 \end{pmatrix},
\qquad
R(s_2) = \begin{pmatrix} 0.3 & -0.5 \\ -0.5 & 1.0 \end{pmatrix}.
\]

\paragraph{Transition kernel $\mathcal{P}(s' \mid s, a_1, a_2)$.}
Successful coordination in the preferred action transitions to the
other state; miscoordination keeps the state unchanged.
\[
\mathcal{P}(\cdot \mid s_1) =
\begin{pmatrix}
  (0.17,\; 0.83) & (0.83,\; 0.17) \\
  (0.83,\; 0.17) & (0.33,\; 0.67)
\end{pmatrix},
\quad
\mathcal{P}(\cdot \mid s_2) =
\begin{pmatrix}
  (0.67,\; 0.33) & (0.17,\; 0.83) \\
  (0.17,\; 0.83) & (0.83,\; 0.17)
\end{pmatrix},
\]
where entry $(a_1, a_2)$ gives $(\mathcal{P}(s_1), \mathcal{P}(s_2))$.

\paragraph{Nash equilibria.}
The state-contingent policy $\pi^*_i(a_1 \mid s_1) = 1$,
$\pi^*_i(a_2 \mid s_2) = 1$ is a strict Nash with
$\mathrm{SW} = 5.0$. The state-blind policies ``always $a_1$'' and
``always $a_2$'' are also strict Nash with
$\mathrm{SW} \approx 3.12$. The $38\%$ welfare gap demonstrates
that state-contingent policies are strictly superior when the
underlying game has multiple states.

\section*{Computational details}

All experiments were implemented in Python~3 using NumPy. The
projected policy gradient update at iteration~$n$ is
$\pi_{n+1} = \mathrm{proj}_\Delta(\pi_n + \gamma_n \hat v_n)$,
where $\mathrm{proj}_\Delta$ is the Euclidean projection onto the
probability simplex and $\hat v_n$ is the exact gradient plus
optional zero-mean Gaussian noise scaled as $\sigma / \sqrt{n}$.
Step-size parameters: $\gamma = 0.3$, $m = 10$, $p = 0.75$
($\gamma_n = 0.3/(n+10)^{0.75}$). Basin-of-attraction heatmaps
use a $35 \times 35$ grid with $2{,}000$ iterations per point.
Geometric convergence plots use $200$ independent trials with
$800$ episodes per restart. Equilibrium selection uses $100$ trials
with $1{,}000$ episodes per restart, up to $K = 15$ restarts. All
experiments use the global random seed~$42$. The stochastic game
experiments use episodic PG with $10$ REINFORCE rollouts per gradient
step; the restart experiment uses $150$ trials with $800$ iterations
per restart. Source code: \texttt{simulations/restart\_pg.py}
(normal-form games) and \texttt{simulations/stochastic\_pg.py}
($|\calS| = 2$ game).

\bibliographystyle{plainnat}
\bibliography{references}

\end{document}
